\medterm{Z-Plasty} A reconstructive operation that rearranges nearby skin in a Z-shaped pattern. It can lengthen a tight scar, change the direction of a scar, and improve movement or appearance.

\medterm{ZAP-70 Deficiency} A rare inherited immune disorder in which T cells cannot develop or work properly. Affected children may have repeated, severe infections and need specialist immune treatment and genetic evaluation.

\medterm{Zafirlukast} A leukotriene-receptor antagonist used for long-term control and prevention of asthma symptoms. It is taken orally and is not a rescue treatment for an acute asthma attack; liver injury and drug interactions are important safety considerations.


\synonyms
Accolate

\medterm{Zanamivir} An inhaled neuraminidase inhibitor used to treat or prevent influenza A and B when started appropriately. It may cause bronchospasm and is used cautiously in people with underlying airway disease.

\synonyms
Relenza

\medterm{Zavanelli Maneuver} An emergency procedure for severe shoulder dystocia that does not respond to other measures. The clinician gently returns the baby’s head to the birth canal, followed by immediate cesarean delivery.

\synonyms
Cephalic Replacement Maneuver

\medterm{Zellweger Syndrome} A severe inherited disorder in which cells cannot make peroxisomes, structures that help process fats and other substances. It begins in infancy and may cause weak muscle tone, feeding problems, seizures, liver disease, and impaired brain development.

\synonyms
ZSD; Cerebrohepatorenal Syndrome

\medterm{Zenker Diverticulum} An acquired, false (pulsion) outpouching of the pharyngeal mucosa and submucosa through Killian's triangle, an area of anatomical muscular weakness between the thyropharyngeus and cricopharyngeus muscles in the posterior hypopharynx. It develops secondary to increased swallowing pressure caused by uncoordinated relaxation or spasm of the cricopharyngeus muscle (upper esophageal sphincter). Most commonly affecting elderly individuals, symptoms include progressive difficulty swallowing (dysphagia), regurgitation of undigested food hours after eating, foul breath (halitosis), a compressible neck lump, and cough or recurrent aspiration pneumonia. Diagnosis is established by a barium swallow study; blind endoscopy or nasogastric tube placement is avoided due to the high risk of puncturing the pouch. Treatment involves cricopharyngeal myotomy with surgical or endoscopic diverticular division.

\synonyms
Zenker's Diverticulum; Pharyngoesophageal Diverticulum; Hypopharyngeal Diverticulum; Pulsion Diverticulum


\medterm{Ziconotide} A strong pain medicine delivered continuously into the fluid around the spinal cord for selected severe chronic pain. It can cause dizziness, confusion, hallucinations, and other serious nervous-system or mental-health effects.

\synonyms
Prialt

\medterm{Zidovudine} An HIV medicine that interferes with the virus’s ability to copy itself. It may be used in selected treatment plans and to reduce transmission around childbirth; important risks include anemia, low white-cell counts, muscle problems, and other toxicity.

\synonyms
Azidothymidine; AZT; Retrovir

\medterm{Ziehl-Neelsen Stain} A laboratory stain used to find bacteria with a waxy coating that resists removal by acid, especially the bacteria that cause tuberculosis. These bacteria appear red against a contrasting background under a microscope.

\synonyms
Acid-Fast Stain; AFB Stain

\medterm{Zieve Syndrome} A temporary complication of heavy alcohol use and sudden liver injury. It causes early breakdown of red blood cells, yellowing of the skin or eyes, and increased blood fats. Treatment includes stopping alcohol and managing liver complications.

\synonyms
Hemolytic Anemia with Hyperlipidemia and Jaundice

\medterm{Zika} An infection caused by Zika virus. It often causes no symptoms or a mild fever, rash, red eyes, joint pain, and muscle pain. Infection during pregnancy can harm the developing baby, and it can rarely trigger Guillain-Barré syndrome.

\synonyms
Zika Virus Disease; ZVD; Congenital Zika Syndrome

\medterm{Zika Virus} A virus spread mainly by infected \textit{Aedes} mosquitoes and also through sexual contact and, less commonly, blood or pregnancy-related transmission. Infection is often mild or symptom-free but can injure a developing fetus.

\synonyms
ZIKV

\medterm{Zileuton} A 5-lipoxygenase inhibitor used for prevention and chronic control of asthma. It does not relieve acute bronchospasm; liver-enzyme elevations and interactions with other medicines require monitoring.

\synonyms
Zyflo

\medterm{Zinc} An essential mineral needed for many body reactions, normal growth, immunity, reproduction, and wound healing. Too little can cause deficiency; excessive supplements can cause stomach symptoms and copper deficiency.

\synonyms
Zn

\medterm{Zinc Acetate} An oral zinc salt used to treat zinc deficiency and, in specific formulations, to reduce intestinal copper absorption in Wilson disease. Dose and monitoring depend on the indication; nausea and gastrointestinal irritation may occur.

\medterm{Zinc Deficiency} A state of inadequate available zinc caused by insufficient intake, malabsorption, excessive losses, or increased requirements. Manifestations may include impaired growth, dermatitis, alopecia, diarrhea, altered taste, delayed wound healing, and increased infection susceptibility.

\synonyms
Hypozincemia; Zinc Depletion

\medterm{Zinc Oxide Ointment} A topical preparation containing zinc oxide that forms a protective barrier over irritated or moist skin. It is commonly used for diaper dermatitis and minor skin irritation.

\synonyms
Topical Zinc Barrier Cream; Zinc Oxide Paste

\medterm{Zinc Oxide} An insoluble zinc compound used topically as a skin protectant and barrier, including in preparations for diaper dermatitis and minor irritation. It can reduce contact with moisture and irritants but does not treat every cause of a rash.

\synonyms
ZnO

\medterm{Zinc Sulfate} An oral or parenteral zinc salt used to prevent or treat zinc deficiency when clinically indicated. Gastrointestinal upset is common with oral therapy, and prolonged excessive use can cause copper deficiency.

\medterm{Zinc Toxicity} Adverse effects caused by excessive zinc ingestion or exposure. Acute toxicity may cause nausea, vomiting, abdominal pain, and diarrhea; chronic excess can impair copper absorption and lead to anemia or neurologic and immune abnormalities.

\synonyms
Hyperzincemia; Zinc Poisoning


\medterm{Zirconia} A zirconium-oxide ceramic used in dentistry for crowns, bridges, implants, and other restorations. Its strength and tissue compatibility make it useful, although restoration design, surface treatment, and clinical indication affect performance.

\synonyms
Zirconium Dioxide; Ceramic Biomaterial


\medterm{Zoledronic Acid} A potent intravenous bisphosphonate used for osteoporosis and certain cancer-related conditions, including hypercalcemia of malignancy and skeletal complications of bone metastases. Renal injury, hypocalcemia, acute-phase reactions, and rare osteonecrosis of the jaw are important risks.

\synonyms
Zoledronate; Zometa; Reclast

\medterm{Zollinger-Ellison Syndrome} A rare condition in which a hormone-producing tumor makes too much gastrin, causing the stomach to produce excess acid. It can lead to repeated ulcers, severe reflux, abdominal pain, and diarrhea.

\medterm{Zolpidem} A nonbenzodiazepine sedative-hypnotic used for short-term treatment of insomnia. It can cause next-day impairment, complex sleep behaviors, dependence, and respiratory or central-nervous-system depression, especially with other sedatives.

\synonyms
Ambien



\medterm{Zonisamide} An antiseizure medicine used mainly for focal seizures. It can cause sleepiness, thinking problems, kidney stones, reduced sweating with overheating, acid imbalance, severe rash, or other serious reactions.

\synonyms
Zonegran


\medterm{Zonule} Fine fibers that hold the eye’s lens in place. They transfer tension from the ciliary muscle so the lens changes shape and the eye can focus at different distances.

\synonyms
Zonule of Zinn; Ciliary Zonule; Suspensory Ligament of Lens

\medterm{Zoon Balanitis} A chronic benign inflammation of the glans penis, usually seen in uncircumcised men. It commonly appears as a sharply demarcated shiny orange-red or rust-colored patch; diagnosis is clinical, with biopsy used when the diagnosis is uncertain.

\synonyms
Plasma Cell Balanitis; Balanitis Circumscripta Plasmacellularis

\medterm{Zoonosis} An infection that naturally spreads between animals with backbones and people. Examples include rabies, brucellosis, leptospirosis, toxoplasmosis, and some influenza infections; spread may involve bites, food, water, droplets, or direct contact.

\synonyms
Zoonotic Infection; Zoonotic Disease

\medterm{Zoster} A painful, usually unilateral vesicular eruption caused by reactivation of latent varicella-zoster virus in a sensory nerve ganglion. It may be complicated by postherpetic neuralgia, ocular disease, or neurologic involvement.

\synonyms
Herpes Zoster; Shingles; Zona

\medterm{Zygomatic Bone} A paired facial bone that forms the prominence of the cheek and part of the lateral orbital wall and floor. It articulates with the frontal, sphenoid, temporal, and maxillary bones.

\synonyms
Cheekbone; Malar Bone

\medterm{Zygomatic Fracture} A break in the cheekbone or the group of bones joining it to the eye socket, usually caused by facial trauma. It may flatten the cheek, numb the skin below the eye, cause double vision, limit mouth opening, or injure the eye socket.

\synonyms
ZMC Fracture; Malar Fracture; Tripod Fracture

\medterm{Zygote} The first cell formed when a sperm joins an egg. It contains genetic material from both parents, divides repeatedly, and develops into an embryo.

\synonyms
Fertilized Oocyte; Conceptus